\chapter{Criterio de cierre operativo de Tomo III}

\section{Cierre operativo}
\begin{definition}
Se llama cierre operativo de Tomo III al estado disciplinado en el que la capa metodológica, curatorial y de proyección de claims alcanza una forma suficientemente completa, trazable y no contradictoria como para declararse estabilizada en su función.
\end{definition}

\begin{observation}
El cierre operativo no equivale a clausura absoluta de todo el sistema MADELEINE. Su función es fijar cuándo Tomo III ya cumple de manera suficiente su papel dentro de la arquitectura general del corpus.
\end{observation}

\section{Condición mínima de cierre}
\begin{definition}
Existe cierre operativo de Tomo III cuando pueden declararse simultáneamente:
\begin{enumerate}
  \item una admisión metodológica legible;
  \item una curación suficiente de demos y claims;
  \item una continuidad explícita entre demo, claim, auditoría, prioridad y salida pública;
  \item una arquitectura externa sobria compatible con el estado alcanzado del corpus;
  \item ausencia de contradicción inmediata entre la función interna del tomo y su proyección externa.
\end{enumerate}
\end{definition}

\begin{observation}
No se exige todavía agotamiento de todas las posibles demos, claims o expansiones del sistema. Basta con que Tomo III ya no dependa de vacíos estructurales que impidan reconocer su función.
\end{observation}

\section{Cierre y relación con Tomo I y Tomo II}
\begin{definition}
Todo cierre operativo de Tomo III debe preservar la diferencia funcional entre:
\begin{enumerate}
  \item Tomo I como núcleo y canon basal;
  \item Tomo II como canon intermedio y organización operativa;
  \item Tomo III como capa metodológica, curatorial y de proyección disciplinada.
\end{enumerate}
\end{definition}

\begin{observation}
Esta cláusula impide que el cierre operativo de Tomo III pretenda refundar el fundamento o absorber las funciones de los tomos anteriores.
\end{observation}

\section{Cierre y reservas}
\begin{definition}
El cierre operativo de Tomo III es compatible con la existencia de material observado, claims reservados y futuras expansiones, siempre que tales reservas no destruyan la legibilidad de su función actual.
\end{definition}

\begin{observation}
La presencia de reservas no invalida el cierre operativo. Lo invalida solo si vuelve indiscernible la arquitectura ya alcanzada del tomo.
\end{observation}

\section{Condición mínima de no cierre}
\begin{definition}
No debe declararse cierre operativo de Tomo III cuando la relación entre demos, claims, auditoría, prioridad y proyección pública siga dependiendo de lagunas estructurales, ambigüedad funcional o colapso entre niveles del corpus.
\end{definition}

\begin{observation}
La no declaración de cierre operativo no implica fracaso del tomo. Puede indicar solo que todavía conviene consolidar mediaciones internas antes de fijar su forma estable.
\end{observation}

\section{Cláusula de alcance}
\begin{observation}
El criterio de cierre operativo de Tomo III fija la primera condición de estabilización funcional del nivel metodológico de MADELEINE. No constituye todavía una clausura absoluta del sistema ni una prohibición de futuras expansiones disciplinadas del corpus.
\end{observation}

